\chapter{Neo4j Database}

Neo4j is a graph database efficient at modeling relationships of entities. It is not a semantic web triplestore, yet they share some similarities and have overlaps in use cases.

In this appendix chapter, Neo4j is briefly introduced. More details are given in Neo4j website \cite{neo4j2026}. 

\section{Introduction}

A brief introduction of Neo4j is given. A comparison of it with semantic web is given.

\subsection{Neo4j as a Graph Database}

TBA

\subsection{AuraDB}

Neo4j can be deployed either as a fully-managed cloud database, in which case, it is also known as \mync{AuraDB}, or a self-managed local database. 

The fully-managed AuraDB introduces additional running cost. In exchange, it is often more robust than self-managed databases and it comes with additional features such as data importer and analyzer.

The self-managed local database has a community version which is fully functional and free-of-charge, and an enterprise version which has all the features in the community version plus addons such as clustering.

In the remaining of this chapter, self-managed community version of Neo4j database is used in all demonstrations.

\subsection{Semantic Web VS Neo4j}

Notice that Neo4j is not a semantic web triplestore, as some of the features of semantic web, such as hierarchical class, properties and OWL are not (fully) supported, and SPARQL is not supported. As a result, it cannot provide as much description logic reasoning as a standard semantic web benchmarked by the late SPARQL standards.

Both being graphic databases, Neo4j and a typical semantic web store data in different manners. The underlying data structures are different, hence the different supported query features.

That said, Neo4j is still very useful. It is less complicated than semantic web, and thanks to that it is more efficient and often more user-friendly. It is especially good at graphical pattern matching.

\section{Neo4j Design}

TBA

\section{Neo4j Implementation}

TBA
